\chapter{Research Methodology and System Design}
\label{chap:methodology}

This chapter documents the original construction method and frozen
149-repository v0.0.6 protocol. The subsequent 207-repository v0.0.7 protocol,
including different provider and metric denominators, is specified alongside
its results in Chapter~\ref{chap:v007}. References here to the ``primary''
artifact refer to the original v0.0.6 experiment, not the newest study.

\section{Methodological Position}
This work adopts the six-activity Design Science Research Methodology (DSRM) of
Peffers et al.~\cite{peffers2007dsrm} as its primary reference methodology.
DSRM is appropriate because the central research contribution is an artifact:
a new static-analysis system designed to solve an identified engineering and
evidence problem. Goal--Question--Metric (GQM)~\cite{basili1992gqm} provides the
measurement layer, connecting the research objective to explicit questions
and reproducible metrics. The empirical component is a frozen
multi-repository observational study; it evaluates artifact behaviour without
turning static findings into unverified human-defect labels.

The unit of execution is a repository fixed at a recorded commit SHA. The unit
of source selection is a classified manifest row. The unit of provider output
is a raw finding. The principal unit of analysis is a canonical occurrence.
Keeping these units separate prevents repository counts, files, raw provider
records and canonical evidence from being mixed in one denominator.

\subsection{Research Design at a Glance}
The methodological choices are deliberately hierarchical rather than a list of
unrelated methods. Table~\ref{tab:research_design_glance} states the role of
each element.

\begin{table}[H]
\centering
\caption{Research design and the role of each methodological element.}
\label{tab:research_design_glance}
\small
\begin{tabular}{p{3.2cm}p{4.0cm}p{7.0cm}}
\toprule
Layer & Choice & Function in this project \\
\midrule
Research paradigm & Design science & Creates and evaluates AnsibleGuard as an
artifact that addresses a demonstrated review-evidence problem. \\
Primary process & Peffers' DSRM & Organises problem identification, objectives,
design, demonstration, evaluation, and communication. \\
Measurement framework & GQM & Converts the research goal into RQ1--RQ4,
metrics, evidence sources, and claim boundaries. \\
Engineering process & Iterative, evidence-driven development & Uses focused
fixtures, real failures, regression tests, and versioned contracts to refine
the artifact between releases. \\
Empirical method & Repository mining with a frozen observational corpus &
Evaluates v0.0.6 on 149 immutable revisions without changing the projects. \\
Change-attribution method & Paired controlled comparison & Compares v0.0.5 and
v0.0.6 on identical repositories, files and LOC to remove population drift. \\
Primary artifact & AnsibleGuard v0.0.6 & The system whose coverage, execution
health, outputs, and behaviour are evaluated. \\
Primary unit of analysis & Canonical occurrence & A source-linked evidence unit
after cross-provider normalisation, with provenance retained. \\
\bottomrule
\end{tabular}
\end{table}

DSRM is the methodology of the research; it is not a claim that the software
was implemented in six linear passes. The implementation itself was iterative.
GQM is not a second competing methodology: it is the measurement instrument
inside DSRM evaluation. Likewise, the 149-repository scan is not the complete
methodology; it is one empirical evaluation within Activity~5. This separation
keeps the report reproducible and prevents a development workflow from being
misreported as a research method.

\subsection{Instantiation of the Six DSRM Activities}
\begin{figure}[H]
\centering
\resizebox{\textwidth}{!}{%
\begin{tikzpicture}[
 d/.style={draw=ansblue,rounded corners,fill=anslightblue,text width=28mm,
 minimum height=12mm,align=center,font=\small\bfseries},
 a/.style={-{Latex[length=2mm]},thick,ansblue},node distance=6mm]
\node[d] (p) {1. Problem identification};
\node[d,right=of p] (o) {2. Objectives of a solution};
\node[d,right=of o] (b) {3. Design and development};
\node[d,right=of b] (demo) {4. Demonstration};
\node[d,below=of demo] (eval) {5. Evaluation};
\node[d,left=of eval] (comm) {6. Communication};
\draw[a] (p)--(o); \draw[a] (o)--(b); \draw[a] (b)--(demo);
\draw[a] (demo)--(eval); \draw[a] (eval)--(comm);
\draw[a,dashed] (eval) to[bend left=25] (b);
\end{tikzpicture}}
\caption{DSRM instantiated for the AnsibleGuard project. The dashed feedback
edge represents evidence-driven detector and semantic-engine refinement.}
\label{fig:dsrm_instantiation}
\end{figure}

\begin{description}
  \item[Problem identification.] Ansible review evidence is fragmented across
  structural checks, runtime-sensitive variables, external linters, secrets,
  Git review metadata and incompatible outputs. Failures can be confused with
  clean results, and raw totals can be mistaken for defects.
  \item[Objectives.] Create one auditable system that covers the review
  checklist where automation is defensible, preserves provider health and
  provenance, models repository-level semantics, supports PR/commit/local
  analysis, and exports a stable validated contract.
  \item[Design and development.] Build the parser boundary, artifact model,
  detector catalogue, symbol/execution graph, provider adapters,
  canonicalisation/comparison engine, profiles and output-v3 package.
  \item[Demonstration.] Execute focused fixtures, real repositories, pull
  requests and the frozen 149-repository population at immutable SHAs.
  \item[Evaluation.] Apply unit/integration/contract/package gates, provider
  health checks, arithmetic reconciliation, paired same-population comparison,
  distribution analysis and manual mechanism sampling.
  \item[Communication.] Publish source, schemas, manifests, commands, study
  summaries, limitations, this report and the supervisor presentation.
\end{description}

The evaluation addresses four research questions:
\begin{description}
  \item[RQ1 -- Coverage.] Which checklist concerns and Ansible quality
    mechanisms are operationalised?
  \item[RQ2 -- Execution reliability.] Can every frozen repository be analysed
    without converting provider or parser failures into clean results?
  \item[RQ3 -- Behavioural change.] On an identical population, how do raw
    findings, canonical occurrences, and semantic states change by version?
  \item[RQ4 -- Diagnostic mechanisms.] Which runtime-variable mechanisms
    dominate the remaining \code{unguarded-variable} cases?
\end{description}

\subsection{Goal--Question--Metric Traceability}
\begin{longtable}{p{1.2cm}p{4.5cm}p{8.2cm}}
\toprule RQ & Evaluation goal & Metrics and evidence \\ \midrule\endhead
RQ1 & Determine implemented and bounded review coverage. & Checklist mapping
(15 Yes, 18 Partial, 5 No), 76-concept catalogue, categories, profiles,
provider ownership and non-automatable boundaries. \\
RQ2 & Determine whether execution evidence remains trustworthy. & Repository
completion; selected files/LOC; provider completed/partial/failed/unavailable;
diagnostics by stage/status; schema validation and arithmetic invariants. \\
RQ3 & Explain artifact evolution without population drift. & Same 149 SHAs,
27,709 files and 1,178,908 LOC; v0.0.5/v0.0.6 raw, canonical, uncertain and
non-uncertain counts; mechanism-linked code/test evidence. \\
RQ4 & Identify remaining variable-provenance limitations. & Unguarded-variable
status totals; unresolved repository concentration; definition/guard/runtime/
order/uncertainty mechanisms; targeted source samples. \\
\bottomrule
\end{longtable}

Metrics are accepted only when their population, source, transformation and
interpretation are stated. Counts are descriptive. A ratio is not called
precision unless its denominator is an independently labelled set of reported
occurrences; recall additionally requires known relevant defects.

Precision and recall are outside the confirmed results because no complete,
independently adjudicated v0.0.6 label set exists. A future protocol is
specified in Section~\ref{sec:future_eval}; it is not reported as completed.

\subsection{Requirements-to-Evidence Traceability}
The artifact was not considered successful merely because it executed. The
following acceptance requirements connect the original problem to an
implementation mechanism and to observable evaluation evidence.

\begin{longtable}{@{}p{0.8cm}p{3.1cm}p{4.8cm}p{5.0cm}@{}}
\caption{Artifact requirements and acceptance evidence.}
\label{tab:requirements_traceability}\\
\toprule ID & Requirement & Realisation & Acceptance evidence \\
\midrule\endfirsthead
\toprule ID & Requirement & Realisation & Acceptance evidence \\
\midrule\endhead
R1 & Tolerant source ingestion & YAML/Jinja boundary, artifact classification,
and coordinate preservation & Parser fixtures; invalid/encrypted files remain
diagnostics; selected-file and LOC reconciliation. \\
R2 & Review-scope coverage & Checklist mapping, 76 concepts, 12 catalogue
categories, and explicit manual boundaries & Audited 15 Yes, 18 Partial, 5 No
mapping and complete rule catalogue. \\
R3 & Repository semantics without global leakage & Static execution graph,
persistent-definition propagation, ordered facts, lexical locals, and scoped
guards & Cross-file include/import tests plus disconnected-file and Jinja
\code{for}/\code{endfor} counterexamples. \\
R4 & Specialist-tool collaboration & Isolated adapters for ansible-lint and
GitLeaks with project-root planning & Provider fixtures; nested-root collapse;
independent configured roots remain distinct. \\
R5 & No false clean result on execution failure & Provider health and parser
failures are diagnostics independent from smells & 149/149 repository
completion with 21 partial ansible-lint runs and 200 visible diagnostics. \\
R6 & Auditable evidence & Canonical occurrences retain every contributing
provider and source coordinate & Raw/canonical reconciliation, provider
evidence arrays, fingerprints, and source-coordinate sampling. \\
R7 & Safe change attribution & File-aware occurrence identity and explicit
legacy comparison path & Tests reject cross-file generic location matches;
paired runs require identical population identity. \\
R8 & Portable, validated outputs & Shared runtime validator and packaged
canonical output-v3 schema copied into scan directories & Schema/contract tests
and wheel-install validation outside the checkout. \\
R9 & Fit-for-purpose use & Profiles and local, commit, PR, and repository modes
& Profile inventory, command workflows, and user/output decision mapping. \\
R10 & Reproducible evaluation & Frozen SHAs, manifests, configurations,
provider versions, summaries, and validation gates & Same 149 repositories,
27,709 files, and 1,178,908 LOC in the paired comparator. \\
\bottomrule
\end{longtable}

\section{Artifact Construction from First Principles}
AnsibleGuard was developed as a new Python package rather than a wrapper around
one linter. Construction proceeded by: (1) translating 38 review criteria into
automatable, partial, and non-automatable concerns; (2) defining stable rule,
category, severity, status, and evidence contracts; (3) implementing tolerant
YAML/Jinja parsing and artifact classification; (4) adding repository symbol
and execution graphs; (5) integrating ansible-lint and GitLeaks as specialist
providers; (6) canonicalising overlapping evidence; and (7) verifying the
package, schema, comparisons, and frozen-corpus behaviour.

The audited checklist coverage is 15 \emph{Yes}, 18 \emph{Partial}, and 5
\emph{No}. These values describe functional scope, not accuracy.

\section{Repository Population Selection and Freezing}
\label{sec:repository_selection}
The 149 repositories were not selected by convenience and were not ranked by
stars. Selection Protocol v2 defines the target population as public,
non-fork, non-archived, non-disabled, non-mirror Ansible repositories whose
latest commit was within six years of the frozen cutoff, with at least 1,000
commits, more than 50 pull requests, at least one confirmed production Ansible
file, and at least 50 merged pull requests changing production Ansible.
Therefore, the study represents mature, active, public repositories suitable
for both snapshot and PR analysis; it does not represent every Ansible project.

\begin{figure}[H]
\centering
\resizebox{\textwidth}{!}{%
\begin{tikzpicture}[
 b/.style={draw=ansblue,rounded corners,fill=anslightblue,minimum height=10mm,
 text width=25mm,align=center,font=\scriptsize\bfseries},
 a/.style={-{Latex[length=2mm]},thick,ansblue},node distance=5mm]
\node[b] (discover) {Inclusive GitHub + Galaxy discovery};
\node[b,right=of discover] (dedup) {Canonical repository deduplication};
\node[b,right=of dedup] (cheap) {Metadata, recency, status and commit gates};
\node[b,right=of cheap] (clone) {Clone and freeze immutable SHA};
\node[b,right=of clone] (class) {Classify every YAML candidate};
\node[b,below=of class] (prs) {Inspect merged PR changed paths};
\node[b,left=of prs] (final) {Apply production-Ansible and PR gates};
\node[b,left=of final] (manifest) {Write final frozen manifest};
\node[b,left=of manifest] (scan) {Scan the 149 eligible revisions};
\draw[a] (discover)--(dedup); \draw[a] (dedup)--(cheap);
\draw[a] (cheap)--(clone); \draw[a] (clone)--(class);
\draw[a] (class)--(prs); \draw[a] (prs)--(final);
\draw[a] (final)--(manifest); \draw[a] (manifest)--(scan);
\end{tikzpicture}}
\caption{Repository discovery, screening, freezing, and study inclusion.}
\label{fig:repository_selection}
\end{figure}

Discovery is intentionally inclusive. GitHub queries include forks in the raw
ledger so the later fork gate can report attrition; Galaxy role and collection
registries contribute source repositories; canonical owner/repository
identities are deduplicated. The complete 2026-08-05 GitHub discovery contains
667,044 ledger rows. Strong topic/Galaxy evidence is enriched before weak
metadata-only evidence for cost control, but both strata remain recorded.

Each gate is tri-state (true, false, or not evaluated), and no candidate row is
deleted. The sequential gate order is: uniqueness; public visibility; not
disabled; not forked; not archived; not mirrored; recent activity; at least
1,000 commits; more than 50 total PRs; confirmed production Ansible; and at
least 50 merged Ansible-changing PRs. Stars remain descriptive only. Applying
the inexpensive commit gate before PR-file inspection avoids costly API work
for ineligible candidates without changing the rule itself.

Repositories passing early gates are cloned. The default branch is resolved to
an immutable SHA and verified with the local Git history. Every YAML candidate
then receives a row-level classification and reason. Production playbooks,
tasks, handlers, defaults, vars, inventory and metadata are separated from CI,
tests, documentation examples, fixtures, generated/vendor content, historical
Ansible, invalid YAML, encrypted vault files and unclear YAML. The manifest
hash, frozen SHA and classifier outcome are retained with the scan.

Finally, merged PRs are checked using exact changed paths and rename/removal
evidence. A PR qualifies only when it changes production Ansible; YAML in CI,
Molecule, tests, documentation, examples or release notes does not qualify by
extension alone. The resulting final-frozen manifest supplies exactly the 149
repository paths and SHAs used by every paired run. This prevents population
drift between v0.0.5 and v0.0.6.

\subsection{Exact Eligibility Rules for the 149-Repository Population}
Repository selection rules and source-analysis rules serve different purposes.
The former define which projects may enter the study; they do not indicate
whether a repository is ``good'' or ``bad.'' A candidate enters the final
population only when every mandatory gate in
Table~\ref{tab:corpus_gates} is true. Unknown values remain
\code{not-evaluated} until enrichment and cannot silently pass.

\begin{longtable}{@{}p{0.8cm}p{4.0cm}p{4.8cm}p{4.0cm}@{}}
\caption{Sequential Selection Protocol v2 eligibility gates.}
\label{tab:corpus_gates}\\
\toprule Step & Gate & Operational rule & Purpose \\
\midrule\endfirsthead
\toprule Step & Gate & Operational rule & Purpose \\
\midrule\endhead
1 & Unique repository & One canonical owner/name identity after source
deduplication & Prevent duplicate sampling across discovery channels. \\
2 & Public & Repository is publicly accessible & Permit reproducible inspection
and redistribution of metadata. \\
3 & Not disabled & GitHub disabled flag is false & Exclude inaccessible or
administratively disabled projects. \\
4 & Not a fork & Fork flag is false & Preserve repository independence and
avoid duplicate upstream histories. \\
5 & Not archived & Archived flag is false & Retain actively maintained targets. \\
6 & Not a mirror & Mirror URL is absent & Avoid duplicated development history. \\
7 & Recent activity & Latest commit falls within six years of the frozen cutoff
& Bound inactive-project bias while retaining mature infrastructure projects. \\
8 & Development history & At least 1,000 commits & Select projects with enough
history for mature snapshot and review analysis. \\
9 & Review activity & More than 50 total pull requests & Ensure an observable
review process. \\
10 & Production Ansible & Repository classification confirms production
Ansible artifacts & Exclude projects where YAML exists only in CI, tests,
examples, fixtures, or unrelated configuration. \\
11 & Ansible review history & At least 50 merged PRs change classified
production Ansible paths & Ensure sufficient Ansible-specific review history. \\
\bottomrule
\end{longtable}

Stars, popularity bands, repository size, and descriptive project type are
retained for analysis but are not eligibility gates. The final attrition step
reduced 265 PR-audited candidates to 149 after the requirement of at least 50
merged production-Ansible-changing PRs; every exclusion remains recorded in
the ledger rather than being deleted.

\section{Evidence-Preserving Pipeline}
\begin{figure}[H]
\centering
\begin{tikzpicture}[
 box/.style={draw=ansblue,rounded corners,fill=anslightblue,
 minimum height=10mm,text width=25mm,align=center,font=\small\bfseries},
 arr/.style={-{Latex[length=2.5mm]},thick,ansblue},node distance=6mm and 7mm]
\node[box] (d) {Repository discovery};
\node[box,right=of d] (p) {YAML/Jinja parsing};
\node[box,right=of p] (c) {Artifact classification};
\node[box,right=of c] (g) {Execution graph};
\node[box,below=of g] (r) {Rules and providers};
\node[box,left=of r] (k) {Canonical merge};
\node[box,left=of k] (v) {Output-v3 validation};
\node[box,left=of v] (s) {Reports and summaries};
\draw[arr] (d)--(p); \draw[arr] (p)--(c); \draw[arr] (c)--(g);
\draw[arr] (g)--(r); \draw[arr] (r)--(k); \draw[arr] (k)--(v);
\draw[arr] (v)--(s);
\end{tikzpicture}
\caption{Analysis pipeline. Failures remain diagnostics and never become
zero-finding results.}
\label{fig:method_pipeline}
\end{figure}

Files are selected from a manifest and parsed with tolerant YAML handling.
Jinja expressions are extracted separately. Each file is classified as a
playbook, tasks, handler, defaults, vars, metadata, inventory, group variables,
host variables, or import-only artifact. Rules run only on compatible classes.
Parse failures retain their source and diagnostic type.

\subsection{End-to-End Execution of the Frozen Study}
The exact operational sequence is shown in
Table~\ref{tab:study_execution_pipeline}. It begins before scanning and ends
only after aggregate and contract validation.

\begin{longtable}{@{}p{0.7cm}p{3.3cm}p{4.9cm}p{4.9cm}@{}}
\caption{Numbered pipeline used for the 149-repository v0.0.6 study.}
\label{tab:study_execution_pipeline}\\
\toprule No. & Name & Input and transformation & Persisted evidence \\
\midrule\endfirsthead
\toprule No. & Name & Input and transformation & Persisted evidence \\
\midrule\endhead
1 & Discover & Inclusive GitHub topic/search and Galaxy source discovery;
canonical deduplication & Complete candidate ledger and discovery provenance. \\
2 & Screen & Apply cheap metadata, status, recency, commit, and PR gates in
sequence & Tri-state gate columns, attrition summaries, exclusion reasons. \\
3 & Freeze & Clone eligible candidates, resolve default branch, record and
verify immutable commit SHA & Repository path, URL, SHA, commit date. \\
4 & Classify & Inspect every YAML candidate and separate production artifacts
from CI, tests, documentation, examples, vendor files, invalid YAML, and
encrypted content & Row-level file
manifest, artifact type, selection reason, LOC. \\
5 & Audit PRs & Inspect exact changed paths of merged PRs against the production
classification & Qualifying PR count, scan progress, final eligibility. \\
6 & Plan analysis & Load the final manifest and \code{technical-full} profile;
resolve artifact compatibility and ansible-lint project roots & Population
identity, profile/configuration, execution plan. \\
7 & Native analysis & Parse selected files, construct repository symbol and
execution graphs, and run compatible AnsibleGuard rules & Native raw findings,
source regions, status and uncertainty evidence. \\
8 & External providers & Run isolated ansible-lint roots and GitLeaks, retaining
partial/error/unavailable states & Raw provider findings, versions, health, and
diagnostics. \\
9 & Canonicalise & Normalise concepts, merge overlapping evidence, compute
file-aware fingerprints, preserve all contributors & Canonical occurrence
array plus provider-evidence arrays. \\
10 & Validate/export & Enforce output-v3 schema and runtime types, copy the
canonical schema, write reports and per-repository summaries & JSON, CSV,
SARIF/text/Markdown projections, schema and metadata. \\
11 & Reconcile & Aggregate 149 outputs, check counts, locations, population
identity, completion, provider health, and paired-version invariants &
All-repository summary and validation evidence used in Chapter~\ref{chap:results}. \\
\bottomrule
\end{longtable}

\section{Execution and Variable Semantics}
Let $F$ be the selected file set and $E_s\subseteq F\times F$ the statically
resolved \code{include_tasks}/\code{import_tasks} edges. A persistent runtime
definition created in $f_i$ is visible in $f_j$ only if a directed path exists
through $E_s$ and program order permits the definition to precede the use.
There is no repository-wide union of facts, so disconnected files do not
inherit definitions. Dynamic include paths contribute named uncertainty.

For a variable use $u(v)$, the detector evaluates
\[
S(u,v)=\langle D_{static},D_{runtime},G_{guard},U_{external},O_{order}\rangle,
\]
recording static definitions, persistent runtime producers, dominating guards,
external/dynamic uncertainty, and ordering evidence. A state is
\code{UNRESOLVED} only when no definition, guard, producer, or uncertainty
source is visible. A use before a later producer is \code{ORDERED_CONFLICT}.
Jinja \code{for} variables are scoped to the matching \code{for}/\code{endfor}
region.

\section{Rule and Provider Responsibilities}
The catalogue contains 76 concepts across correctness, reliability,
maintainability, style, compatibility, security, documentation, testing,
structure, dependencies, and delegated lint concerns. Native AnsibleGuard rules
own semantic and repository-level predicates. ansible-lint owns established
Ansible conventions such as FQCN, naming, YAML style, and module practices.
GitLeaks supplies candidate-secret pattern evidence.

Provider health is orthogonal to source quality. \code{completed},
\code{partial}, and \code{tool-error} are diagnostics. Missing executables,
timeouts, and crashes are never source smells and never imply a clean scan.
Nested unconfigured ansible-lint roots collapse to prevent double linting;
independently configured roots remain separate.

\section{Canonical Occurrences and Output v3}
Provider findings are normalised into occurrences containing rule/concept,
category, severity, semantic status, artifact, file/line, message, confidence,
uncertainty reasons, fingerprint, and contributing provider evidence.
Canonicalisation merges overlap without erasing provenance. Severity is the
maximum contributor under \code{note < warning < error}.

Generic location-only findings cannot match across files as semantic renames.
Comparison requires file-aware identity or stronger semantic evidence. Output
v3 is a versioned envelope checked by a shared validator. The canonical JSON
Schema is packaged in the wheel and copied into each scan output. Legacy lists
remain supported only through an explicit compatibility path.

\section{Frozen-Corpus Study Protocol}
\label{sec:study_design_m}
The population comprises 149 repositories at recorded SHAs. The manifest
universe contains 64,876 YAML files and 4,254,017 lines; selection yields
27,709 Ansible files and 1,178,908 lines. v0.0.5 and v0.0.6 are compared only
when repository, commit, files, and LOC are equal, eliminating population drift
as an explanation.

Each repository records completion, files/LOC, raw provider findings, canonical
occurrences, statuses, severities, categories, concepts, artifacts, and
diagnostics. Aggregate arithmetic is reconciled with per-repository outputs.
Coordinates are checked against existing files and line ranges. Suspicious
concentrations are sampled by mechanism rather than inferred from totals.

\subsection{Operational Measures and Computation Rules}
Let $R$ be the frozen repository set, $F_m$ the manifest YAML files, $F_s$ the
selected Ansible files, $L_s$ the selected LOC, $H_p$ the repositories on which
provider $p$ completed, $N_{raw}$ the raw provider findings, and $N_c$ the
canonical occurrences. The principal descriptive measures are:
\begin{align}
\text{selection rate} &= \frac{|F_s|}{|F_m|}\times100, \\
\text{provider completion}_p &= \frac{|H_p|}{|R|}\times100, \\
\text{canonicalisation reduction} &=
  \frac{N_{raw}-N_c}{N_{raw}}\times100, \\
\text{canonical density} &= \frac{N_c}{L_s}\times1000, \\
\text{status share}_s &= \frac{N_{c,s}}{N_c}\times100.
\end{align}
For v0.0.6 these yield a 42.71\% file-selection rate, a 6.22\%
raw-to-canonical reduction, and 261.22 canonical occurrences per selected
KLOC. Provider completion is always reported together with partial,
tool-error, and unavailable states. Density is a review-workload normaliser,
not a defect prevalence estimate.

\begin{longtable}{@{}p{3.2cm}p{4.1cm}p{2.9cm}p{3.8cm}@{}}
\caption{Data collection and interpretation plan.}
\label{tab:data_collection_plan}\\
\toprule Measure & Authoritative source & Aggregation & Permitted interpretation \\
\midrule\endfirsthead
\toprule Measure & Authoritative source & Aggregation & Permitted interpretation \\
\midrule\endhead
Repository identity & Frozen manifest and resolved Git SHA & Exact equality,
then count & Evaluation population and repeatability. \\
Files and LOC & Row-level selection manifest & Sum only selected rows & Analysis
denominator; never inferred from finding files. \\
Raw findings & Per-provider result arrays & Sum before canonicalisation &
Provider workload, including overlap. \\
Canonical occurrences & Validated output-v3 occurrence arrays & Sum after
canonicalisation & Auditable evidence workload, not human-confirmed defects. \\
Semantic status & Occurrence \code{status} & Count by enumerated state & Degree
and kind of semantic evidence or uncertainty. \\
Provider health & Provider evidence and diagnostics & Count repositories by
state & Execution completeness independent of source quality. \\
Version delta & Paired summaries with population invariant & Subtract only
after identity checks & Behavioural change attributable to artifact version. \\
Mechanism sample & Occurrence, source context, and detector trace & Group by
mechanism and inspect exemplars & Hypothesis generation and targeted diagnosis,
not population precision. \\
\bottomrule
\end{longtable}

\section{Verification Ladder}
Evidence is accepted after focused regression tests, full pytest, documentation
staleness checks, Python compilation, diff validation, wheel installation
outside the checkout, schema tests, available remote-equivalence and benchmark
gates, paired frozen-corpus comparison, and manual mechanism sampling.

\section{Validity and Claim Boundaries}
The design establishes reproducible static evidence and explicit execution
health on the frozen population. It does not establish ecosystem prevalence,
universal recall, or a human true-positive label for every occurrence.
\code{CONFIRMED} means a detector predicate was satisfied without a specialised
uncertainty state; it is not manual adjudication.

\section{Future Labelled Evaluation}
\label{sec:future_eval}
A future precision study should freeze version and corpus, publish guidance for
true-positive, false-positive, and borderline labels per rule, draw a
reproducible stratified sample, obtain two
independent human labels, report Cohen's $\kappa$, adjudicate disagreements,
and publish weighted Wilson intervals. Recall requires a separate mutation or
expert-curated benchmark. Until those artifacts exist, this report uses
descriptive and diagnostic terminology rather than precision claims.
